\documentclass[11pt,a4paper]{article}
\usepackage[T1]{fontenc}
\usepackage{lmodern,amsmath,amssymb,amsthm,mathtools,microtype,booktabs}
\usepackage[margin=26mm]{geometry}
\usepackage[colorlinks=true,linkcolor=blue!45!black,citecolor=blue!45!black,urlcolor=blue!45!black]{hyperref}
\newtheorem{theorem}{Theorem}[section]
\newtheorem{proposition}[theorem]{Proposition}
\newtheorem{lemma}[theorem]{Lemma}
\newtheorem{corollary}[theorem]{Corollary}
\theoremstyle{definition}
\newtheorem{definition}[theorem]{Definition}
\newtheorem{remark}[theorem]{Remark}
\newcommand{\Q}{\mathbb Q}
\newcommand{\Z}{\mathbb Z}
\newcommand{\OK}{\mathcal O_K}
\newcommand{\OL}{\mathcal O_L}
\newcommand{\Tr}{\operatorname{Tr}}
\newcommand{\Nm}{\operatorname{N}}
\newcommand{\A}{\mathbb A}
\setlength{\emergencystretch}{2em}
\title{Existential Definitions of Rings of Integers and Subfields}
\author{}
\date{9 November 2026}
\begin{document}
\maketitle
\begin{abstract}
For a number field $K$, the established Diophantine definability of
$\Z$ in $\mathcal O_K$ makes the definability of $\mathcal O_K$ in
$K$ equivalent to that of $\Z$ in $K$. Assuming an existential
definition of $\Z$ in $\Q$, these properties are also equivalent to
existential definability of the rational subfield in $K$.
We give explicit bounds on coordinates and degree for these
reductions and for descent through a finite field extension.
For every finite extension $L/K$ of characteristic-zero fields,
existential definability of $K$ in $L$ is equivalent to definability
of the relative norm image, the norm graph, or the trace graph.
The norm-image implication follows from an exact finite-difference
identity. In contrast, the norm and trace graphs of a proper
extension are Zariski dense and have no quantifier-free definitions.
Finally, a fixed finite collection of existential formulas cannot
define rings of integers throughout all finite extensions of a
number field. These results distinguish reductions and obstructions
to transferring a formula from a counterexample to the existence of
a separately chosen formula over each number field.
\end{abstract}

\section{Definitions and arithmetic scope}

A subset $S\subseteq F^r$ of a field $F$ is \emph{Diophantine in $F$}
if it is the projection onto the first $r$ coordinates of the common
zero set of finitely many polynomials with coefficients in $F$.
Thus there are polynomials $f_1,\ldots,f_s$ such that
\[
 \mathbf x\in S\quad\Longleftrightarrow\quad
 \exists\mathbf y\in F^m:\quad
 f_1(\mathbf x,\mathbf y)=\cdots=f_s(\mathbf x,\mathbf y)=0.
\]
Over a field, a nonzero condition $g\ne0$ can be expressed by
$gz=1$. Finite conjunctions, finite unions, and existential
projections of Diophantine conditions are again Diophantine.
For a number field, a finite system can be replaced by one equation;
Lemma~\ref{lem:combine} gives a quantitative version. Coefficients
from the specified ambient field are allowed throughout.

The \emph{total number of coordinates} includes all distinguished
input coordinates as well as all quantified coordinates. Degree is
total polynomial degree. A formula with a single input and $N$
total coordinates has $N-1$ witnesses.

For a number field $K$, write $\OK$ for its ring of algebraic
integers. We use the established theorem
\begin{equation}\label{eq:DL}
            \Z\text{ is Diophantine in }\OK
            \quad\text{for every number field }K.
\end{equation}
The theorem is stated in \cite[Corollary 1.2]{ABHS}; the
rank-stability method and the reduction to all number fields are
explained in \cite[Sections 1.2 and 5]{Poonen}.
All quantified variables in \eqref{eq:DL} range over $\OK$.
This distinction is essential: \eqref{eq:DL} is not itself a
definition of $\OK$ in $K$.

Where needed, we state separately the hypothesis
\begin{equation}\label{eq:ZQ}
                    \Z\text{ is Diophantine in }\Q.
                    \tag{H}
\end{equation}
The elementary implications below do not supply a proof of
\eqref{eq:ZQ}, and their validity does not depend on a particular
proposed polynomial for that hypothesis.

\section{Combining equations and taking linear images}

\begin{lemma}\label{lem:combine}
Over a number field $F$, a system of $s\ge1$ polynomial equations of
degree at most $D$ can be replaced, in the same coordinates, by one
equation of degree at most $\lceil\sqrt{s}\rceil D$.
If $F$ has a real embedding and $s>1$, the bound $2D$ also holds.
For $s=1$, the original equation suffices.
\end{lemma}
\begin{proof}
Put $e=\lceil\sqrt{s}\rceil$. Choose a finite prime $\mathfrak p$
of $F$, an element $\pi$ of valuation one there, and a monic
polynomial $g$ of degree $e$ with integral coefficients and
irreducible reduction over the residue field. Let
\[
 N_g(Z_0,\ldots,Z_{e-1})
  =\det\bigl(m_{Z_0+Z_1u+\cdots+Z_{e-1}u^{e-1}}\bigr)
\]
be the multiplication determinant in $F[u]=F[U]/(g)$.
For any nonzero block $\mathbf z\in F^e$, reduction after scaling
by its least coordinate valuation gives
\[
        v_{\mathfrak p}(N_g(\mathbf z))
                      =e\min_i v_{\mathfrak p}(z_i).
\]
Indeed, the reduced multiplication matrix multiplies by a nonzero
element of the residue extension field, so its determinant is a
unit. Therefore the degree-$e$ form
\[
       A_e(\mathbf z_0,\ldots,\mathbf z_{e-1})
                    =\sum_{j=0}^{e-1}\pi^jN_g(\mathbf z_j)
\]
vanishes only when all $e^2$ entries are zero: its nonzero summands
have distinct valuation classes modulo $e$, so a unique term has
least valuation. Pad the original residuals to length $e^2$ and
substitute them into $A_e$.

If $F$ has a real embedding, a sum of squares vanishes in $F$
only when each summand vanishes. The sum of the squares of the
residuals gives the second bound.
\end{proof}

The following elementary observation avoids an unnecessary
coordinate in several constructions.

\begin{lemma}\label{lem:linear-image}
For $1\le i\le n$, let $S_i\subseteq F$ be defined by a polynomial
$F_i(T,\mathbf Y_i)$ with $N_i$ total coordinates and degree at most
$D$. Fix $b_1,\ldots,b_n\in F$, with $b_n\ne0$. The set
\[
              b_1S_1+\cdots+b_nS_n
\]
has a definition with $\sum_iN_i$ total coordinates, $n$ equations,
and degree at most $D$.
\end{lemma}
\begin{proof}
For the distinguished input $x$, retain $z_1,\ldots,z_{n-1}$ and
put
\[
               z_n=\frac{x-\sum_{i<n}b_i z_i}{b_n}.
\]
Impose $F_i(z_i,\mathbf Y_i)=0$ for every $i$, using disjoint
witness blocks. There are
\[
          1+(n-1)+\sum_i(N_i-1)=\sum_iN_i
\]
coordinates. Substitution of a linear expression does not increase
total degree. The equivalence with membership in the displayed
linear image is immediate in both directions.
\end{proof}

\section{Rational integers, algebraic integers, and rational inputs}

\begin{theorem}\label{thm:equivalence}
For every number field $K$, the following assertions are equivalent:
\begin{enumerate}
\item $\OK$ is Diophantine in $K$;
\item $\Z$ is Diophantine in $K$.
\end{enumerate}
Under hypothesis~\eqref{eq:ZQ}, they are also equivalent to:
\begin{enumerate}\setcounter{enumi}{2}
\item the subfield $\Q$ is Diophantine in $K$.
\end{enumerate}
\end{theorem}
\begin{proof}
Assume (1). Take a definition supplied by \eqref{eq:DL} and replace
the requirement that each input or witness lie in $\OK$ by its
assumed existential definition in $K$. All quantifiers remain
existential. Restriction of every original variable to its original
ring makes the new formula define exactly $\Z$.

Assume (2). Choose an integral basis
$\omega_1,\ldots,\omega_n$ of $\OK$ over $\Z$. Then
\[
 x\in\OK\quad\Longleftrightarrow\quad
 \exists z_1,\ldots,z_n\in\Z:\quad
                    x=\sum_{i=1}^n\omega_i z_i.
\]
Replacing the integer membership conditions by their definitions
inside $K$ proves (1).

Still under (2), rationality is expressed by
\[
 q\in\Q\quad\Longleftrightarrow\quad
 \exists a,b\in\Z\ \exists c\in K:
                         qb=a,\qquad bc=1.
\]
Thus (2) implies (3), without using \eqref{eq:ZQ}.

Finally assume (3) and \eqref{eq:ZQ}. In the defining formula over
$\Q$, restrict its input and every witness to the Diophantine
subfield $\Q\subseteq K$. The resulting existential formula over
$K$ is equivalent to the original one over $\Q$ and defines $\Z$.
\end{proof}

\begin{proposition}\label{prop:integral-count}
Let $n=[K:\Q]$. If $\Z$ has a definition in $K$ with $N$ total
coordinates and degree $D$, then $\OK$ has a definition with
$nN$ total coordinates and $n$ equations of degree at most $D$.
It has a single-equation definition in those coordinates of degree
at most $\lceil\sqrt n\rceil D$, or at most $2D$ if $K$ has a
real embedding and $n>1$.
\end{proposition}
\begin{proof}
Apply Lemma~\ref{lem:linear-image} to the integral-basis identity
in the proof of Theorem~\ref{thm:equivalence}, then apply
Lemma~\ref{lem:combine}.
\end{proof}

\begin{proposition}\label{prop:composition-count}
Suppose that a polynomial of degree $D$ and $N$ total coordinates
defines $\Z$ in $\Q$, and a polynomial of degree $E$ and $R$ total
coordinates defines $\Q$ in a number field $K$.
Then $\Z$ has a definition inside $K$ with $NR$ total coordinates,
$N+1$ equations, and degree at most $\max(D,E)$.
The ring $\OK$ has a definition with $nNR$ total coordinates,
$n(N+1)$ equations, and the same degree bound, where $n=[K:\Q]$.
\end{proposition}
\begin{proof}
Write the first polynomial as $P(T_0,\ldots,T_{N-1})$ and the
second as $Q(T,\mathbf V)$, with $R-1$ witnesses in $\mathbf V$.
Impose
\[
 P(T_0,\ldots,T_{N-1})=0,\qquad
 Q(T_i,\mathbf V_i)=0\quad(0\le i<N).
\]
The $N$ original coordinates and the $N(R-1)$ new ones total $NR$.
The argument in Theorem~\ref{thm:equivalence} proves correctness.
For $\OK$, use $n$ copies of this system in the integral-basis
construction and eliminate the last integer coordinate as in
Lemma~\ref{lem:linear-image}. This yields the stated count without
first aggregating either system.
\end{proof}

\begin{remark}
In the first part of Proposition~\ref{prop:composition-count},
restricting only the witnesses is insufficient in general: a
polynomial with rational coefficients and rational values for the
witnesses can still have nonrational algebraic roots in its input.
In particular, the polynomial defining $\Z$ over $\Q$ cannot simply
be evaluated with unrestricted $K$-valued witnesses.
\end{remark}

\section{Relative norms and traces characterize the missing subfield}

Throughout this section, $L/K$ is a finite extension of
characteristic-zero fields of degree $n$. The norm and trace are
the usual determinant and trace of multiplication on the
$n$-dimensional $K$-vector space $L$. Define the following subsets
of powers of the \emph{ambient field $L$}:
\[
\begin{split}
 T_{L/K}&=\{\Nm_{L/K}(a):a\in L\}\subseteq K\subseteq L,\\
 G_N&=\{(a,b)\in L^2:b=\Nm_{L/K}(a)\},\\
 G_T&=\{(a,b)\in L^2:b=\Tr_{L/K}(a)\}.
\end{split}
\]
In particular, definability here does not silently require the
coordinates of $a$ in a $K$-basis to belong to $K$.

\begin{lemma}\label{lem:finite-difference}
For every integer $n\ge1$, one has the polynomial identity
\begin{equation}\label{eq:difference}
 \sum_{j=0}^{n-1}(-1)^{n-1-j}\binom{n-1}{j}(T+j)^n
                    =n!T+\frac{n!(n-1)}2.
\end{equation}
If $S\subseteq K$ contains every $n$th power of an element of $K$,
then
\begin{equation}\label{eq:linear-image}
       K=\sum_{j=0}^{n-1}
                   (-1)^{n-1-j}\binom{n-1}{j}S.
\end{equation}
\end{lemma}
\begin{proof}
The left side of \eqref{eq:difference} is the $(n-1)$st forward
difference of $T^n$. It is linear with coefficient $n!$.
Its constant term equals $(n-1)!S(n,n-1)$, where $S(n,n-1)$
counts partitions of an $n$-element set into $n-1$ nonempty parts.
Exactly one pair is joined, so $S(n,n-1)=\binom n2$, giving the
stated constant. This also follows directly by inclusion and
exclusion applied to surjections from an $n$-element set to an
$(n-1)$-element set. For $n=1$, the identity is $T=T$.

Every element on the right of \eqref{eq:linear-image} belongs to
$K$. Conversely, given $q\in K$, choose
$T=q/n!-(n-1)/2\in K$ in \eqref{eq:difference}. Each $(T+j)^n$
belongs to $S$, and their indicated linear combination is $q$.
\end{proof}

\begin{theorem}\label{thm:norm-trace}
The following assertions are equivalent:
\begin{enumerate}
\item $K$ is Diophantine in $L$;
\item the norm image $T_{L/K}$ is Diophantine in $L$;
\item the norm graph $G_N$ is Diophantine in $L$;
\item the trace graph $G_T$ is Diophantine in $L$.
\end{enumerate}
If $T_{L/K}$ is defined by a polynomial with $M$ total coordinates
and degree $D$, then $K$ has a definition with $nM$ total
coordinates and $n$ equations of degree at most $D$.
\end{theorem}
\begin{proof}
Assume (1), and choose a $K$-basis $b_1,\ldots,b_n$ of $L$.
For $a\in L$, write $a=\sum_i a_i b_i$ with all $a_i\in K$.
The entries of its multiplication matrix are fixed $K$-linear
forms in the $a_i$. Its determinant and trace are therefore
polynomials over $K$ in these coordinates. Restricting each
$a_i$ to the given Diophantine subfield $K$ makes the graph of
each map Diophantine in $L$. Projection of the norm graph proves
(2) as well as (3) and (4).

Assertion (3) implies (2) by projection. If (2) holds, then
$T_{L/K}\subseteq K$ and
$\Nm_{L/K}(q)=q^n$ for every $q\in K$.
Lemma~\ref{lem:finite-difference} expresses all of $K$ as a fixed
linear image of $n$ copies of $T_{L/K}$, proving (1).
Lemma~\ref{lem:linear-image} proves the quantitative assertion,
since the last coefficient in \eqref{eq:linear-image} is one.

Finally the trace is onto $K$: for $q\in K$,
$\Tr_{L/K}(q/n)=q$. Projection of $G_T$ therefore proves that
(4) implies (1).
\end{proof}

\begin{corollary}\label{cor:quadratic}
If $L/K$ is quadratic, existential definability in $L$ of any of
the four sets in Theorem~\ref{thm:norm-trace} is also equivalent
to existential definability of the graph of the nontrivial
$K$-automorphism $\sigma$ of $L$.
Moreover,
\[
 x\in K\quad\Longleftrightarrow\quad\sigma(x)=x
 \quad\Longleftrightarrow\quad\Tr_{L/K}(x)=2x
 \quad\Longleftrightarrow\quad\Nm_{L/K}(x)=x^2.
\]
\end{corollary}
\begin{proof}
A definition of the automorphism graph defines its fixed field.
Conversely, if $K$ is Diophantine in $L$, write
$x=a+b\theta$ with $a,b\in K$ and use
$\sigma(x)=a+b\sigma(\theta)$.
The trace equivalence follows from
$\Tr(x)=x+\sigma(x)$.
For the norm, $\Nm(x)=x\sigma(x)=x^2$ implies either $x=0$ or
$\sigma(x)=x$; zero already belongs to $K$.
\end{proof}

For example, when $n=2$, the norm-image identity is simply
\[
                         K=T_{L/K}-T_{L/K}.
\]
Indeed, for $q\in K$, the identity
$((q+1)/2)^2-((q-1)/2)^2=q$ exhibits the difference of two norms
of elements of $K$ itself.

\subsection{Why an ordinary norm form does not solve this problem}

A norm form with \emph{$K$-valued coordinates} is an ordinary
polynomial over $K$. The map on a \emph{single $L$-valued input}
addressed in Theorem~\ref{thm:norm-trace} is a different object.
The distinction has the following exact algebraic manifestation.

\begin{proposition}\label{prop:dense}
If $L/K$ is proper, the graphs $G_N$ and $G_T$ are Zariski dense
in $\A^2_L$. Neither has a quantifier-free definition in the
language of rings with coefficients in $L$. In particular neither
is the graph of a polynomial function in one variable over $L$.
\end{proposition}
\begin{proof}
Choose $\theta\in L\setminus K$. The set $K^2$ is Zariski dense
in $\A^2_L$, because a polynomial vanishing on $K^2$ is zero by
applying the one-variable root bound successively.

Consider first the trace graph. Its points include the images of
$K^2$ under the $L$-linear map
\[
    (a,b)\longmapsto
        \bigl(a+b\theta,\;na+b\Tr_{L/K}(\theta)\bigr).
\]
Its determinant is $\Tr_{L/K}(\theta)-n\theta\ne0$:
otherwise $\theta\in K$. Hence this map is an automorphism of
$\A^2_L$, proving density of $G_T$.

For the norm, restrict the polynomial norm form to the two
coordinates $a,b$ corresponding to $a+b\theta$ and consider
\[
        (a,b)\longmapsto
          \bigl(a+b\theta,\;\Nm_{L/K}(a+b\theta)\bigr).
\]
At $(1,0)$ its Jacobian determinant is again
$\Tr_{L/K}(\theta)-n\theta\ne0$, since the linear term in
$\Nm(1+b\theta)$ is $b\Tr(\theta)$. The map is dominant.
Its values on $K^2$ belong to $G_N$, so $G_N$ is dense as well.

A quantifier-free formula defines a constructible set. A dense
constructible subset of $\A^2_L$ contains a nonempty Zariski
open subset. Since $L$ is infinite, such an open subset has
infinitely many points in at least one vertical line; for instance,
choose a first coordinate at which a nonzero defining polynomial
of a principal open subset remains nonzero. A graph of a function
has exactly one point in each vertical line. This contradiction
rules out a quantifier-free definition. A polynomial graph is a
proper closed subset and is excluded in particular.
\end{proof}

\section{Descent and the structure of a possible counterexample}

\begin{theorem}\label{thm:descent}
Let $L/K$ be a finite extension of number fields of degree $d$.
If $\OL$ has a definition in $L$ by a polynomial with $N$ total
coordinates and degree $D$, then $\OK$ has a definition in $K$
with $1+d(N-1)$ total coordinates, $d$ equations, and degree at
most $D$. A single equation in those coordinates can be chosen
of degree at most $\lceil\sqrt d\rceil D$; if $K$ has a real
embedding and $d>1$, the bound $2D$ also applies.
\end{theorem}
\begin{proof}
Fix a $K$-basis $b_1,\ldots,b_d$ of $L$ and write the defining
polynomial as $P(T,Y_1,\ldots,Y_{N-1})$.
Keep the input $x\in K$ and substitute
\[
                     Y_j=\sum_{i=1}^d b_i y_{ji},
                     \qquad y_{ji}\in K.
\]
Expanding in the basis gives a unique expression
\[
 P\left(x,\sum_i b_i y_{1i},\ldots,\sum_i b_i y_{N-1,i}\right)
                 =\sum_{i=1}^d b_i P_i(x,(y_{ji})),
\]
where every $P_i$ belongs to $K[x,(y_{ji})]$ and has degree at
most $D$. Vanishing is equivalent to all $P_i=0$ by linear
independence. The substitutions range over all original witnesses
in $L$, so the input projection is $\OL\cap K=\OK$.
The coordinate count and aggregation bounds follow immediately.
\end{proof}

\begin{corollary}\label{cor:counterexample}
If some number field $K$ has a ring of integers that is not
Diophantine in $K$, then every finite extension of $K$ has the
same property. In particular, there exists a totally imaginary
number field, Galois over $\Q$, with that property.
\end{corollary}
\begin{proof}
The first assertion is the contrapositive of
Theorem~\ref{thm:descent}. Take a normal closure of $K/\Q$
and adjoin $i$. The resulting field is finite and Galois over
$\Q$, contains $K$, and has no real embedding.
\end{proof}

This is a reduction of a possible counterexample, not the
construction of one. In particular, a polynomial that fails after
base extension does not establish the hypothesis of
Corollary~\ref{cor:counterexample}.

\section{A finite collection of unchanged formulas is insufficient}

An existential formula in this section may use finite Boolean
combinations of polynomial equalities, so nonzero conditions and
finite alternatives are allowed. Coefficients are fixed in a
number field $k$. Its interpretation in an extension $L/k$ uses
the same formula and the same coefficients, with witnesses in $L$.

\begin{lemma}\label{lem:generic}
Let $\Phi(T)$ be such a formula and suppose $\Phi(k)$ is infinite.
There are an integer $d\ge1$ and a nonzero $h\in k[T]$ such that,
for every $a\in k$ with $h(a)\ne0$, the formula $\Phi(a)$ holds
in some extension $E_a/k$ of degree at most $d$.
\end{lemma}
\begin{proof}
The constructible set of witness tuples has a constructible image
in the input line by Chevalley's theorem \cite{Stacks}.
An infinite image is dense in that line, so one conjunction in
the formula has nonempty generic fiber over $k(T)$.
Choose a closed point of this generic fiber, with residue field
$E$ finite over $k(T)$. Write $E=k(T)[\alpha]$, with monic
minimal polynomial $G(T,Z)$ of degree $d$ in $Z$.
Express the witness coordinates and the inverses of all required
nonzero expressions as polynomials in $\alpha$ with coefficients
in $k(T)$.

After inverting a nonzero $h\in k[T]$, all these expressions and
the coefficients of $G$ belong to $R=k[T,1/h]$. The algebra
$C=R[Z]/(G)$ is free of rank $d$ and embeds in $E$.
All equations and inverse identities therefore hold in $C$.
For $h(a)\ne0$, specialize to $C_a=k[Z]/(G(a,Z))$ and pass to
any residue field at a maximal ideal. Its degree over $k$ is
at most $d$, and all equations and nonzero conditions survive.
\end{proof}

\begin{theorem}\label{thm:finite-menu}
Let $\Phi_1,\ldots,\Phi_r$ be a finite nonempty list of
existential formulas over a number field $k$. Suppose
$\Phi_i(k)$ is infinite for every $i$.
There exists $B\ge1$ such that, for all but finitely many
$q\in\Q\setminus\Z$, there is a finite Galois extension
$L_q/k$ of degree at most $B$ in which every $\Phi_i(q)$ holds.
Thus none of the formulas defines $\mathcal O_{L_q}$ in $L_q$.
One may further require $L_q$ to have no real embedding, at the
cost of at most doubling $B$.
\end{theorem}
\begin{proof}
Apply Lemma~\ref{lem:generic} to each formula, obtaining $d_i$
and $h_i$. Exclude the finite union of the rational zero sets
of the $h_i$. For every remaining rational noninteger $q$,
choose the witnessing fields $E_{i,q}$.
The normal closure of $E_{i,q}/k$ has degree at most $d_i!$,
by its permutation action on the conjugates of a primitive
element. Take $L_q$ to be the compositum of these normal closures.
Then $L_q/k$ is Galois of degree at most
\[
                              B=\prod_{i=1}^r d_i!.
\]
Each witness tuple remains a witness in $L_q$: polynomial
equalities and nonzero conditions are preserved by field
extensions. A rational algebraic integer is an integer, so
$q\notin\mathcal O_{L_q}$. Adjoining $i$ gives the last assertion.
\end{proof}

\begin{corollary}
A fixed finite collection of formulas over $k$, each correct for
$\mathcal O_k$ in $k$, cannot supply a definition of the ring of
integers in every finite extension by choosing one formula from
that collection for each extension.
\end{corollary}
\begin{proof}
Each formula has infinitely many inputs over $k$. The extension
in Theorem~\ref{thm:finite-menu} makes every formula in the
collection incorrect simultaneously.
\end{proof}

For a single nonzero polynomial $P(T,\mathbf Y)$ that depends on
a witness coordinate $Y_j$, the integer $d$ in
Lemma~\ref{lem:generic} may be chosen as $\deg_{Y_j}P$.
Indeed, specialize all other witnesses to constants in $k$
so that the leading coefficient in $Y_j$ remains a nonzero
polynomial of $T$. Outside its finite zero set, adjoining a
root gives the asserted extension. Consequently the factorial
bound in Theorem~\ref{thm:finite-menu} can use the smallest
positive individual witness degree of each polynomial.

\section{What rank stability supplies inside a field}

A positive-rank elliptic curve whose rank does not grow in an
extension supplies a useful Diophantine subset of the smaller
field. It is necessary to distinguish this conclusion from
existential definability of the entire smaller field.

\begin{proposition}\label{prop:stable-set}
Let $K$ be a number field and let $E/\Q$ be an elliptic curve
such that
\[
                0<\operatorname{rank}E(\Q)
                         =\operatorname{rank}E(K).
\]
There is an integer $r\ge1$ for which the set
\[
 T=\{x(P):P\in rE(K),\ P\ne O\}
\]
is an infinite Diophantine subset of $K$ contained in $\Q$.
\end{proposition}
\begin{proof}
The Mordell--Weil groups are finitely generated and have the same
rank, so $E(K)/E(\Q)$ is finite. Choose $r$ annihilating this
quotient. Then $rE(K)\subseteq E(\Q)$, and $rE(K)$ is infinite.
The graph of multiplication by the fixed integer $r$ is
algebraic. On a finite collection of affine charts it is given
by polynomial equations and nonzero conditions. Projecting its
nonidentity image onto the abscissa coordinate gives a
Diophantine set over $K$. The abscissa map has finite fibers,
so this set is infinite; all its elements are rational because
its points lie in $E(\Q)$.
\end{proof}

The known arithmetic transfer theorem under the same rank
hypothesis defines the smaller \emph{ring of integers} inside
the larger ring of integers; the more general abelian-variety
version is \cite[Theorem 4.2]{Poonen}. For fields themselves,
\cite[Theorem 4]{CornelissenShlapentokh} gives a definition of
the smaller subfield using one universal quantifier. Remark~6.3 of~\cite{CornelissenShlapentokh} expressly distinguishes
that conclusion from an existential definition in the whole
field. Those results therefore do not remove the rational-input
restriction in Proposition~\ref{prop:composition-count}.

In general, taking multiplicative combinations of the abscissas
in Proposition~\ref{prop:stable-set} does not yield every rational
number. The following obstruction gives an explicit example.

\begin{lemma}\label{lem:abscissa-parity}
Let
\[
                       E:y^2=x^3+Ax+B
\]
be a nonsingular elliptic curve over $\Q$. Suppose that $p$ is
an odd prime, $A$ and $B$ are $p$-integral, and the reduction
of $B$ is a nonsquare in $\mathbb F_p^\times$.
Every affine rational point of $E$ has nonzero abscissa $x$
and even valuation $v_p(x)$. The multiplicative subgroup of
$\Q^\times$ generated by all these abscissas consists entirely
of numbers of even $p$-adic valuation.
\end{lemma}
\begin{proof}
If $x=0$, then $y^2=B$, contrary to the nonsquare residue of
$B$. If $v_p(x)>0$, the equation has a unit right-hand side
congruent to $B$ modulo $p$, giving the same contradiction.
If $v_p(x)<0$, then
\[
        3v_p(x)<v_p(Ax),\qquad 3v_p(x)<v_p(B).
\]
The term $x^3$ has the unique least valuation on the right,
so $2v_p(y)=3v_p(x)$. This forces $v_p(x)$ to be even.
The remaining case has valuation zero. Valuations add under
multiplication and change sign under inversion, proving the
last assertion.
\end{proof}

\begin{proposition}\label{prop:abscissa-example}
The curve $E:y^2=x^3-2$ has a rational point of infinite order,
but $5$ is not a ratio of two rational abscissas of $E$, or
any finite product of integral powers of its rational abscissas.
\end{proposition}
\begin{proof}
The point $(3,5)$ lies on $E$. The curve has good reduction at
$5$ and $7$, since its discriminant is $-1728$. Direct counting
gives
\[
                       \#E(\mathbb F_5)=6,
                       \#E(\mathbb F_7)=7.
\]
For completeness, modulo $5$ the numbers of ordinates above
$x=0,1,2,3,4$ are $0,2,2,1,0$. Modulo $7$ the corresponding
numbers are $0,0,0,2,0,2,2$. Add the identity in each case.

Prime-to-$p$ rational torsion injects under good reduction at
$p$. One way to see the needed injectivity is to use the formal
parameter $z=-x/y$ in the reduction kernel: if $p\nmid m$,
then $[m]z=mz+O(z^2)$ has the same finite valuation as $z$ and
cannot vanish. Thus a prime divisor $\ell\ne5,7$ of a rational
torsion order would divide both $6$ and $7$, which is impossible.
The $5$-primary torsion is excluded by reduction at $7$, and
the $7$-primary torsion by reduction at $5$. Hence $E(\Q)$
has trivial torsion, and $(3,5)\ne O$ has infinite order.

At $p=5$, the constant coefficient is $B=-2\equiv3$, a
nonsquare. Lemma~\ref{lem:abscissa-parity} shows that every
multiplicative combination in the statement has even
$5$-adic valuation, whereas $v_5(5)=1$.
\end{proof}

The obstruction in Proposition~\ref{prop:abscissa-example}
concerns multiplicative combinations. It does not apply to
sums of such values and does not disprove a more elaborate
existential definition. In particular, it supplies neither a
counterexample to Theorem~\ref{thm:equivalence}'s equivalent
properties nor a positive result for a new number field.

\section{Consequences for constructions over number fields}

The distinction between two questions is exact. Theorems
\ref{thm:descent} and \ref{thm:equivalence} concern the existence
of a formula chosen for a specified field. Theorem
\ref{thm:finite-menu} concerns reusing a fixed finite collection
of formulas throughout extensions. The second question has a
negative answer even if the first has a positive answer for
every number field.

Under \eqref{eq:ZQ}, a proof that $\OK$ is Diophantine in $K$
can therefore proceed by constructing any one of the following:
a definition of the rational subfield; a definition of the
absolute norm image as a subset of $K$; or a definition of the
absolute trace graph with witnesses in $K$.
Theorem~\ref{thm:norm-trace} shows that these are equivalent
arithmetic tasks, rather than independent shortcuts obtained
by writing a formal norm or trace symbol.
Alternatively, a direct definition of any subset $S\subseteq\OK$
containing $\Z$ is sufficient: the integral-basis construction
gives
\[
                    \OK=\sum_{i=1}^n\omega_i S,
\]
because the right side is contained in $\OK$ and contains its
full integral-basis expansion. This last sufficient condition
does not assume \eqref{eq:ZQ}.

No theorem above gives an example of a number field whose ring
of integers is not Diophantine in its fraction field, or a proof
that all number fields have the property. The remaining step is
an arithmetic definition inside the field itself, with all
witnesses in the field and with soundness against nonintegral
inputs.

\begin{thebibliography}{9}
\bibitem{ABHS}
L.~Alp\"oge, M.~Bhargava, W.~Ho and A.~Shnidman,
\emph{Rank stability in quadratic extensions and Hilbert's tenth
problem for the ring of integers of a number field},
Invent. Math. \textbf{243} (2026), 1129--1139.
\url{https://doi.org/10.1007/s00222-025-01392-3}.

\bibitem{Poonen}
B.~Poonen, \emph{Rank stability makes rings of integers diophantine},
3 July 2026.
\url{https://math.mit.edu/~poonen/papers/h10_over_OK.pdf}.

\bibitem{CornelissenShlapentokh}
G.~Cornelissen and A.~Shlapentokh,
\emph{Defining the integers in large rings of a number field using
one universal quantifier},
J. Math. Sci. (N.Y.) \textbf{158} (2009), 713--726.
\url{https://doi.org/10.1007/s10958-009-9404-4};
preprint \url{https://arxiv.org/abs/0708.3075v2}.

\bibitem{Stacks}
The Stacks Project Authors, \emph{The Stacks Project},
Chevalley's theorem, Tag 00FE.
\url{https://stacks.math.columbia.edu/tag/00FE}.
\end{thebibliography}
\end{document}
